
RAG~\citep{lewis2020retrieval, gao2023retrieval} addresses fundamental limitations of specific knowledge coverage in LLMs by integrating external knowledge retrieval with generative models~\citep{wang2023augmenting}.
%RAG improves factual consistency and enables access to real-time information, effectively mitigating the problems of knowledge timeliness and hallucinated generation. 
%However, recent advances that enhance RAG capabilities through modularity, autonomy, and coverage mechanisms~\citep{singh2025agentic} typically introduce substantial computational overhead through billion-parameter LLMs, limiting their practical deployment in real-life settings.



\pb{Query Enhancement Techniques.}
Query is one of the factors that determine the success of retrieval mechanism for the generator~\cite{wang2024evaluating}.
Thus, several works explore the RAG improvement technique from the query enhancement perspective.
RQ-RAG~\citep{chan2024rq} enables query rewriting, decomposition, and disambiguation. 
DMQR-RAG~\citep{li2024dmqr} proposes diverse multi-query rewriting strategies operating at different information levels. 
HyDE~\citep{gao2022precise}  generates hypothetical documents to bridge the semantic gap between queries and relevant passages. 
Step-back Prompting~\citep{zheng2023take} rewrites detailed queries at higher conceptual levels. 
In addition, recent multimodal reasoning studies emphasize structured intermediate reasoning processes (e.g., Visual Chain-of-Thought) to enhance alignment between queries and retrieved evidence, further highlighting the importance of query abstraction and decomposition in complex reasoning tasks~\citep{zheng2026towards}.
Our method is inspired by these works, thus leveraging a query rewriting mechanism in the pre-retrieval process through both decomposition and enhancement.
%While methods enhance retrieval diversity by exploring multiple query perspectives, they typically operate as pre-retrieval stages without systematic coverage detection.


\pb{Evidence Retrieval Methods.}
Contemporary RAG systems are generally grouped into two based on retrieval mechanisms, sparse and dense.
Sparse retrieval refers to utilizing keywords to look for information in the database with BM25 or TF-IDF mechanisms~\cite{ram2023context, huly2024old}.
Meanwhile, dense retrieval leverages similarity measures between the semantic representation of the user's query and the candidate evidence~\cite{borgeaud2022improving, izacard2022few, lewis2020retrieval}.
Additionally, recent works combine sparse and dense retrievals. 
RAG-Fusion~\citep{rackauckas2024rag} applies reciprocal rank fusion to combine results from multiple query reformulations.
DAT~\citep{hsu2025dat} proposes dynamic alpha tuning for hybrid retrieval, adapting weight coefficients based on individual query characteristics.
Late-interaction retrieval models such as ColBERTv2 improve retrieval granularity by preserving token-level interactions, enabling more precise matching between queries and passages~\citep{santhanam2022colbertv2}.
Beyond these paradigms, recent efforts explore more fine-grained and efficient retrieval mechanisms. For instance, HASH-RAG introduces deep hashing techniques to enable efficient yet precise retrieval~\citep{guo2025hash}, while instruction-driven overlap reduction further improves retrieval efficiency by minimizing redundant evidence~\citep{ou2025accelerating}.
%rather than using fixed weights across all queries. 
%Beyond weight tuning, instruction-based fine-tuning enables lightweight models to effectively perform both sparse and dense retrieval~\cite{lin2023ra}.
Nevertheless, these approaches assume that the post-retrieval condition is perfect and thus, the evidence is directly passed to the generator.
%These approaches ensure both exact entity matching and semantic generalization in the retrieval process.


\pb{Evidence Integrity Checking.} 
Post-retrieval mechanisms, including reranking and context compression, address evidence quality.
RankRAG~\citep{yu2024rankrag} unifies context ranking with retrieval-augmented generation, using LLMs to rank retrieved passages before feeding them to the generator. 
Late Chunking~\citep{gunther2024late} addresses semantic fragmentation by encoding documents with full in-document context before chunking, preserving contextual information that fixed-length chunking typically loses. 
Structure-aware approaches~\citep{zhang2026text} further tackle this by incorporating global document structure into the evidence processing pipeline. 
Despite the growing research in semantic chunking, the proposed methods generally incur high computational cost, which cannot be sufficiently justified by the gains proposed~\citep{qu2025semantic}.
%These methods, however, incur a high computational cost from embedding the full document in a high-dimensional embedding space.
%Recent work identifies semantic fragmentation from fixed-length document chunking as a critical problem affecting both retrieval and generation quality; most existing solutions focus on improving chunking strategies at the indexing stage rather than addressing truncation post-retrieval.


%\pb{Autonomous and Self-Correcting Systems.} 
%Self-RAG~\citep{asai2024self}  enables mid-generation self-reflection and dynamic retrieval decisions. 
%CRAG~\citep{yan2024corrective} introduces lightweight evaluators to assess retrieval quality and trigger corrective web searches when documents fail to support generation. 
%These systems implement feedback loops for correction, typically operating at the generation or retrieval evaluation stages rather than throughout the entire pipeline.


%\paragraph{Agentic RAG Systems.} 



\pb{Computationally Efficient RAG.} 
To address the increasingly complex RAG framework and the utilization of large generative models as a retriever, current studies explore ways to make LLM and RAG more efficient~\citep{zhou2026look,wang2026stream,wang2026efficient}.
MiniRAG~\citep{fan2025minirag} utilizes heterogeneous graph indexing, combining text chunks with entity hierarchies to reduce reliance on dense embeddings. 
Plan\textsuperscript{*}RAG~\citep{verma2024plan} structures multi-hop reasoning as directed acyclic graphs (DAGs), enabling systematic exploration of reasoning paths with parallel execution for efficiency. 
LightRAG~\citep{guo2024lightrag} incorporates graph structures into text indexing with dual-level retrieval, enhancing contextual awareness while improving response times.
Atlas~\citep{izacard2022atlas} demonstrates that jointly pretraining retrieval and generation at scale can significantly improve knowledge-intensive task performance while maintaining efficiency.
Recent works also investigate model compression and lightweight architectures for multimodal and language models, such as knowledge restructuring techniques that transform external information into structured triplets for more efficient retrieval~\citep{wang2026transforming} and
lightweight agent memory mechanisms built upon small language models~\citep{zhang2026lightweight}.
Furthermore, advances in efficient multimodal agent learning highlight the role of transferable experience in reducing redundant computation during reasoning~\citep{li2026experience}. 
Our work is inspired by these studies in leveraging smaller-sized yet powerful retriever to reduce the computational cost of RAG framework.
%These approaches demonstrate that architectural design can reduce computational requirements.%, though they primarily optimize data structures and retrieval routing rather than embedding systematic semantic reasoning throughout the pipeline.